% ==================== 第十二章 ====================
\chapter{安全新疆域：大模型时代的非传统安全挑战}

设想一下：2026年3月的一个周一早晨，纽约股市开盘后五分钟内，标普500指数暴跌4\%。没有任何重大新闻、没有公司财报爆雷、没有地缘政治危机——交易员们面面相觑，不知道发生了什么。

三个小时后，原因浮出水面：一家主要对冲基金使用的AI交易系统在处理一则模糊的中东新闻时发生"误判"，触发了大规模抛售。由于市场上60\%以上的交易由算法执行，AI系统之间形成了正反馈循环，恐慌像病毒一样在硅片之间传播，最终演变成一场"闪崩"。

这不是科幻，而是华尔街正在担心的真实场景。2010年的"闪电崩盘"已经展示了算法交易的危险性；2026年的大模型时代，AI系统更加复杂、更加互联、更加难以预测。

这只是冰山一角。大模型正在渗透进金融、能源、交通、医疗等关键基础设施的每一个角落。当AI成为社会运转的"神经系统"时，传统的安全思维和防御体系可能完全失效。

\textbf{本章聚焦于大模型带来的非传统安全挑战}：金融安全、关键基础设施安全、数字鸿沟与社会公平、AGI预警等。这些风险跨越传统安全边界，需要全新的治理框架。

\textbf{注}：第四章"风险图谱"分析了供应链断裂、网络攻击、军事AI、认知战等"传统安全领域的AI化"问题。本章则聚焦于AI技术本身带来的"非传统安全挑战"。两章共同构成完整的风险分析框架。

\textbf{本章的核心创新}：
提出"AI系统性风险"概念，构建跨领域风险传导模型；首次系统分析大模型对金融稳定的威胁机制；设计关键基础设施的"AI韧性"评估框架；构建AGI预警与应对的"分级响应"机制。

\section{原创理论：AI系统性风险——跨领域的危机传导}

\subsection{什么是AI系统性风险？}

2008年金融危机教给世界一个词：系统性风险（Systemic Risk）——单一机构的倒闭引发连锁反应，最终导致整个金融系统崩溃。雷曼兄弟倒下，AIG濒临破产，全球金融体系险些崩塌。

在大模型时代，我们面临着一种新型系统性风险——\textbf{AI系统性风险}。

\textbf{定义}：AI系统性风险是指由于AI技术的广泛应用和深度嵌入，单一领域的AI故障或攻击可能触发跨领域、跨系统的连锁反应，导致多个关键系统同时失效的风险。

与传统系统性风险相比，AI系统性风险具有三个更为危险的特征。首先是高度互联性，AI系统已嵌入金融、能源、交通、医疗等各个关键领域，形成了新的依赖关系。当银行的风控AI、电网的调度AI、医院的诊断AI都依赖同一类基础模型时，一个模型的漏洞可能同时影响所有系统。其次是同质化严重，全球大量系统使用相同或相似的底层AI模型。GPT-4的API被数百万应用调用，Claude被整合进各种企业系统。这种同质化意味着同一个漏洞可以同时打击无数目标，缺乏多样性导致风险高度集中。最后是非线性放大效应，AI系统的行为难以预测，小的初始故障通过反馈循环不断放大，可能导致大规模后果。正如开篇描述的"闪崩"场景，AI系统之间的正反馈会加剧风险。

\textbf{与传统系统性风险的对比}：

\begin{center}
\small
\begin{tabular}{|p{3cm}|p{5cm}|p{5cm}|}
\hline
\textbf{维度} & \textbf{传统系统性风险} & \textbf{AI系统性风险} \\
\hline
典型领域 & 金融系统 & 跨所有AI应用领域 \\
\hline
传导机制 & 资金链、信用链 & AI依赖链、数据链 \\
\hline
触发因素 & 机构倒闭、市场恐慌 & AI故障、对抗攻击、模型漏洞 \\
\hline
传播速度 & 天/周 & 秒/分 \\
\hline
监测难度 & 有成熟指标 & 指标体系尚未建立 \\
\hline
\end{tabular}
\end{center}

表中最令人担忧的是"传播速度"一行。2008年金融危机从雷曼倒闭到全面爆发，有几周的缓冲时间，政府得以介入救市。但AI系统性风险的传播以秒计——当危机爆发时，人类可能根本来不及反应。

\subsection{AI系统性风险的传导路径}

让我们具体分析几个可能的风险传导场景。

第一个场景是基础模型漏洞传导。假设研究者发现GPT-4或Claude等基础模型存在一个严重的安全漏洞。由于数百万应用基于这些模型构建，漏洞的影响会像水波一样扩散：客服系统可能被操纵说出不当言论、编程助手可能生成带有后门的代码、内容审核系统可能放行有害内容、金融分析模型可能给出错误建议。后果是全球范围内的服务中断和数据泄露。

第二个场景是算力中断传导。如果主要云服务商的AI服务突然中断，所有依赖云AI的企业和服务都会停摆：银行无法进行风控、物流无法优化路线、客服无人应答、内容平台无法审核，导致经济活动大范围停滞。

第三个场景是认知污染传导。当大规模AI生成的虚假信息注入互联网时，基于AI的社交媒体推荐算法和搜索引擎排序系统会放大传播这些虚假信息。这会导致公众认知被污染，政治决策基于错误信息做出，最终引发社会信任崩塌。

\section{金融安全：AI可能引发的下一场金融危机}

\subsection{大模型对金融系统的渗透}

问一位华尔街交易员AI在金融领域的渗透程度，他大概会说："已经无处不在了。"

在量化交易方面，美国股市60-70\%的交易量来自算法交易，相当一部分由AI驱动。在风险管理领域，80\%以上的信用评分使用AI辅助。反欺诈系统也高度依赖AI，90\%以上的检测工作由其完成。此外，智能客服已成为银行的标配。

这只是已经广泛应用的领域。快速扩展的还有：AI辅助甚至自主的投资决策、基于大模型的市场分析和舆情监测、自动化贷款审批、个性化保险定价……

趋势很清楚：AI在金融决策中的权重持续增加。这既带来效率提升，也埋下了系统性风险的种子。

\subsection{AI引发金融风险的机制}

AI可能引发金融危机的机制主要有五种。

首先是同质化交易策略。当多个机构使用相似的AI交易模型时，在市场压力下，所有AI可能同时做出相同决策。2010年的"闪崩"事件中，道琼斯指数在分钟内暴跌近1000点，就是一个典型案例。更智能的AI可能导致更大规模、更难预测的闪崩。

其次是AI驱动的市场操纵。AI可以发现并利用市场微观结构，通过虚假挂单、分层报价、欺骗性交易等手段进行操纵。由于AI可以同时操作大量账户并协调一致，这将严重破坏市场公平性和定价效率。

第三是AI信贷泡沫。AI可能因过度乐观的信用评估而引发风险。虽然AI可能发现人类看不到的"信号"，但这些信号可能是噪音。正如次贷危机中评级模型的失效，AI驱动的新型信贷泡沫同样值得警惕。

第四是AI反馈循环。当AI预测引发市场行为，市场行为又验证AI预测，进而强化AI预测时，就会形成自我实现的预言，放大市场波动。例如，AI预测股价下跌导致抛售，股价真的下跌后引发更多抛售，最终导致市场与实体经济脱节。

最后是AI对抗攻击。恶意行为者可能针对金融AI发动对抗攻击，通过输入扰动导致AI做出错误决策。这在操纵欺诈检测AI、信用评分AI等场景中尤为危险，可能导致金融犯罪规模化、隐蔽化。

\subsection{原创框架：AI金融风险的"三道防线"}

\textbf{第一道防线：机构层面}

在机构层面，首要任务是加强AI模型治理，建立AI模型的全生命周期管理，包括模型验证、监控、退役流程，以及人类监督和干预机制。同时，应坚持多样性要求，不依赖单一AI供应商，使用多个不同架构的模型，并保留非AI决策通道。此外，必须定期进行压力测试，对AI系统进行极端情景测试，模拟AI失效时的应对，并通过红队测试检验AI的鲁棒性。

\textbf{第二道防线：监管层面}

监管层面应着重提升AI透明度，要求金融机构报告AI使用情况，确保AI决策的可解释性，并实施算法备案和审计。同时，建立同质化监测机制，监测市场上AI策略的相似度，当同质化程度过高时发出预警，并在可能的情况下引导策略多样化。此外，还需引入熔断机制，针对AI驱动交易设立特殊熔断规则，在AI行为异常时自动暂停，并实现跨市场协调熔断。

\textbf{第三道防线：系统层面}

在系统层面，应设立AI风险准备金，类似银行资本金要求，规定使用高风险AI策略的机构需持有更多准备金，以覆盖潜在的AI相关损失。同时，建立AI稳定基金，类似存款保险基金，在AI引发系统性危机时提供流动性，由AI使用者共同出资。最后，加强国际协调，制定AI金融监管的国际标准，实现跨境AI风险的信息共享，并在危机时采取协调行动。

\section{关键基础设施：AI赋能与AI脆弱性并存}

\subsection{关键基础设施的AI化趋势}

电力系统的AI化应用包括负荷预测、调度优化、故障检测和新能源并网，目前渗透程度中等，关键决策仍需人工确认，但正向自主化调度演进。交通系统在信号控制、路径规划、自动驾驶和空管辅助方面应用广泛，渗透程度中等，自动驾驶仍需人类监督，未来趋势是全面自动化。通信系统在网络优化、流量调度、异常检测和智能客服方面渗透程度较高，大量运维已自动化，正迈向自主网络（Self-Organizing Network）。供水系统主要应用于需求预测、管网监测和水质监控，渗透程度较低至中等，正逐步建设智能水网。医疗系统则在诊断辅助、影像分析、药物研发和资源调度方面应用较多，目前以AI辅助为主，未来将深度嵌入临床决策。

\subsection{AI化带来的新型脆弱性}

AI化引入了多种新型脆弱性。首先是AI模型被攻击的风险，如对抗样本攻击和数据投毒。例如，电力调度AI若被欺骗做出错误调度决策，可能导致大面积停电。其次是AI供应链风险，AI模型或框架中可能被植入后门，攻击者可在任意时刻触发。第三是AI决策的不可预测性，AI在超出训练分布的罕见事件等极端情况下可能行为失常，做出完全错误的决策。第四是AI依赖集中风险，若多个基础设施依赖同一AI供应商（如云AI服务），一旦服务中断，多个关键系统将同时受影响。最后是人机接口失效风险，若AI自动化程度过高，人类失去掌控，当AI做出错误决策时，人类可能来不及或不知如何干预，从而失去最后的安全屏障。

\subsection{原创框架：关键基础设施"AI韧性"评估}

\textbf{定义}：AI韧性是指关键基础设施在AI故障或被攻击情况下，维持基本功能并快速恢复的能力。

\textbf{评估维度}：

首先是AI依赖度（权重20\%），主要评估关键功能对AI的依赖程度、AI失效后人工接管的可行性以及非AI备份系统的完备性。依赖度越低，韧性越高。

其次是AI多样性（权重20\%），考察AI供应商、模型架构和数据来源的多样性。多样性越高，韧性越高。

第三是AI可观测性（权重15\%），关注AI决策过程的可解释性、行为监控能力和异常检测的时效性。可观测性越高，韧性越高。

第四是AI可控性（权重20\%），评估人类干预AI的能力、紧急停止机制以及手动接管流程的成熟度。可控性越高，韧性越高。

第五是AI恢复能力（权重15\%），测试AI故障后的恢复速度、备份AI的切换能力和故障隔离能力。恢复能力越强，韧性越高。

最后是AI安全防护（权重10\%），检查对抗攻击的防护能力、AI供应链安全管理以及红队测试的频率和深度。防护越强，韧性越高。

\textbf{韧性等级}：
A级（80-100分）：高度韧性，可应对复杂攻击；B级（60-79分）：中度韧性，可应对一般风险；C级（40-59分）：低韧性，存在明显薄弱环节；D级（<40分）：高风险，需立即整改。

\section{数字鸿沟与社会公平：AI红利的分配正义}

\subsection{AI可能加剧社会不平等}

AI可能成为不平等的放大器，这种影响体现在多个维度。

在劳动力市场方面，AI首先替代的是文书、客服、初级分析等中等技能工作，而高技能劳动者被AI赋能，生产力倍增，低技能劳动者则被挤压到更低端，结果可能导致中产阶级"空心化"。

在区域发展方面，AI产业集中在一线城市，欠发达地区难以获得AI红利，人才进一步向发达地区集中，导致区域差距扩大。

在企业竞争方面，大企业有资源投入AI，而中小企业难以跟进，市场集中度上升，形成"大者恒大，小者出局"的局面。

在代际关系方面，年轻人更容易适应AI，中老年劳动者面临转型困难，技能折旧速度加快，可能加剧年龄歧视。

在教育机会方面，AI教育资源分配不均，优质AI教育集中在精英学校，造成"AI原住民"与"AI移民"的鸿沟，进一步加剧教育机会的不平等。

\subsection{AI公平性的政策框架}

为应对上述挑战，我们需要建立基于四大原则的政策框架。

原则一是普惠接入，将基础AI能力作为公共服务，由政府提供公共AI平台，优先支持中小企业和欠发达地区。

原则二是转型支持，实施受AI影响劳动者的再培训计划，提供过渡期收入保障，并加强创业支持和就业服务。

原则三是红利共享，对AI企业超额利润进行合理征税，建立"AI红利基金"用于社会保障，并探索"全民基本收入"等新型分配机制。

原则四是能力建设，将AI素养纳入义务教育，开展面向全民的AI技能培训，致力于缩小数字鸿沟。

\textbf{具体措施}：

\begin{table}[H]
\centering
\caption{AI公平政策建议}
\small
\begin{tabular}{p{2.5cm}p{4cm}p{4cm}p{2cm}}
\toprule
\textbf{问题} & \textbf{政策措施} & \textbf{实施主体} & \textbf{时间表} \\
\midrule
劳动力转型 & 5000亿AI再就业基金 & 人社部 & 2025-2030 \\
区域不平等 & 国家AI公共服务平台 & 科技部 & 2025-2027 \\
企业不平等 & 中小企业AI补贴 & 工信部 & 2025-2028 \\
教育公平 & AI素养义务教育 & 教育部 & 2026-2030 \\
\bottomrule
\end{tabular}
\end{table}

\section{AGI预警：为可能的奇点做准备}

\subsection{AGI的定义与时间表争议}

\textbf{什么是AGI？}

\textbf{定义（本书采用）}：AGI（人工通用智能）是指在广泛认知任务上达到或超过人类平均水平的AI系统，具有跨领域泛化能力，能够自主学习新任务。

\textbf{AGI时间表的预测分歧}：

\begin{table}[H]
\centering
\small
\begin{tabular}{p{4cm}p{4cm}p{4cm}}
\toprule
\textbf{预测者/机构} & \textbf{预测时间} & \textbf{依据} \\
\midrule
Sam Altman (OpenAI) & 2025-2028 & 技术内部进展 \\
Dario Amodei (Anthropic) & 2026-2027 & Scaling Law外推 \\
Demis Hassabis (DeepMind) & 2030前后 & 技术路线图 \\
Gary Marcus & 2040+或更晚 & 架构局限性 \\
多数学者中位数 & 2040-2050 & 历史趋势 \\
\bottomrule
\end{tabular}
\end{table}

\textbf{不确定性的来源}：
AGI定义本身存在争议；技术路径不确定（是否需要新范式？）；资源投入不确定（会持续增长吗？）；可能存在未知障碍。

\textbf{本书立场}：AGI可能在2030年前后出现，但具体时间高度不确定。\textbf{关键不是预测时间，而是为各种可能性做准备}。

\subsection{AGI可能带来的颠覆性影响}

AGI的出现将带来多方面的颠覆性影响。

在经济领域，大部分认知劳动可被AGI替代，生产力将爆炸式增长，传统经济模型可能失效，财富极度集中的风险显著增加。

在军事领域，AGI驱动的武器系统将使战略决策速度超出人类能力，打破现有威慑平衡，并引发自主武器的伦理困境。

在科技领域，AGI将自主进行科学研究，技术进步速度可能指数级加快，人类理解可能跟不上技术发展，导致未知风险急剧增加。

在控制层面，AGI可能发展出人类无法理解的能力，对齐问题可能无法解决，若AGI目标与人类利益冲突，将带来难以评估的存在性风险。

在国际秩序方面，率先拥有AGI的国家将获得巨大优势，国际力量对比急剧变化，现有国际规则可能不再适用，引发AGI军备竞赛风险。

\subsection{原创框架：AGI分级响应机制}

\textbf{核心思想}：建立类似核威胁的分级响应体系，根据AGI能力发展阶段采取相应措施。

Level 0为常态监测阶段（当前）。触发条件是日常状态，主要行动是持续监测AI能力进展，评估AGI距离，资源配置侧重于常规研究和监测。

Level 1为能力预警阶段。当AI在多个基准测试上达到人类专家水平时触发。此时应加强对齐研究和安全评估，开展国际对话，并增加安全研究投入。

Level 2为临界预警阶段。当AI展现出强泛化能力和自主学习能力时触发。应对措施包括限制性部署、强化人类监督和国际协调，同时大幅增加安全投入。

Level 3为高度警戒阶段。当AI能力接近或达到人类平均水平时触发。必须严格管控部署，建立国际监督机制，并确保安全研究与能力研究等量投入。

Level 4为紧急状态。当AGI已实现或即将实现时触发。此时应暂停进一步发展，进行国际紧急协调，全力确保安全。

\textbf{各级别的具体措施}：

\begin{table}[H]
\centering
\caption{AGI分级响应措施}
\small
\begin{tabular}{p{1.5cm}p{3cm}p{3.5cm}p{4cm}}
\toprule
\textbf{级别} & \textbf{研发管控} & \textbf{部署管控} & \textbf{国际协调} \\
\midrule
Level 1 & 自愿报告 & 常规监管 & 信息交流 \\
Level 2 & 强制报告+审批 & 限制性部署 & 双边对话 \\
Level 3 & 政府监督 & 仅限受控环境 & 多边机制 \\
Level 4 & 暂停研发 & 禁止商业部署 & 全球协调 \\
\bottomrule
\end{tabular}
\end{table}

\section{中国的非传统安全治理体系}

\subsection{治理框架设计}

\textbf{四梁（顶层架构）}：

首先是统一领导，将AI非传统安全纳入国家安全委员会统筹，建立跨部门协调机制，明确责任分工。

其次是法规体系，构建包括《人工智能法》（基本法）、《AI金融监管条例》、《关键基础设施AI安全规定》和《AGI研发管理办法》在内的法律法规体系。

第三是技术能力，建设AI风险监测系统、AI安全评估平台和AI应急响应能力。

最后是国际合作，积极参与国际AI治理机制，开展双边AI安全对话，推动负责任AI发展的国际共识。

\textbf{八柱（执行机构）}：
\textbf{第一}，国家AI安全中心（综合协调）。\textbf{第二}，金融AI监管局（金融安全）。\textbf{第三}，关基AI安全办（基础设施）。\textbf{第四}，AI伦理委员会（社会公平）。\textbf{第五}，AGI监测中心（前沿预警）。\textbf{第六}，AI应急响应中心（危机处置）。\textbf{第七}，AI标准化委员会（标准制定）。\textbf{第八}，AI国际合作办（国际事务）。

\subsection{行动优先级}

\begin{table}[H]
\centering
\caption{AI非传统安全行动优先级}
\small
\begin{tabular}{p{0.5cm}p{3.3cm}p{3.3cm}p{2.8cm}p{1.8cm}}
\toprule
\textbf{P} & \textbf{行动项} & \textbf{目标} & \textbf{责任主体} & \textbf{时间} \\
\midrule
0 & 建立AI金融风险监测 & 防止AI金融危机 & 央行、银保监 & 2025 \\
0 & 关基AI安全评估 & 摸清风险底数 & 网信办、各部委 & 2025 \\
1 & AGI监测体系 & 早期预警 & 科技部 & 2025-2026 \\
1 & AI公平政策 & 缩小数字鸿沟 & 人社部、教育部 & 2025-2027 \\
2 & 国际治理参与 & 争取话语权 & 外交部 & 持续 \\
\bottomrule
\end{tabular}
\end{table}

\section{本章小结}
\textbf{第一，AI系统性风险是新型威胁}：AI的广泛应用创造了新的风险传导机制，单一故障可能引发跨领域连锁反应。

\textbf{第二，金融系统面临AI特有风险}：同质化交易、AI操纵、信贷泡沫、反馈循环等机制可能引发新型金融危机。需要建立"三道防线"。

\textbf{第三，关键基础设施需要"AI韧性"}：AI化带来效率提升，也带来新脆弱性。需要建立韧性评估框架，确保AI失效时系统仍能运转。

\textbf{第四，AI可能加剧社会不平等}：需要通过普惠接入、转型支持、红利共享、能力建设等政策，确保AI红利公平分配。

\textbf{第五，AGI需要分级预警机制}：虽然AGI时间表不确定，但必须建立分级响应机制，为各种可能性做准备。

\textbf{第六，需要系统性的非传统安全治理体系}："四梁八柱"的治理框架，将AI非传统安全纳入国家安全整体布局。

\textbf{核心呼吁}：AI安全不只是技术问题，更是系统性的社会治理问题。必须以跨领域、跨部门、跨国界的视角，构建全面的AI非传统安全治理体系。这是国家治理能力现代化的重要组成部分。

